\section{Parte 1: problemas abiertos y fronteras del RL profundo}

En esta clase, Chelsea Finn resume los problemas abiertos y las líneas de frontera del RL profundo, y ofrece recomendaciones prácticas para investigar en RL profundo.

\subsection{Desafíos en la definición del problema}

\subsubsection{Diseño de la recompensa}

\begin{importantbox}{La dificultad de fondo del diseño de la recompensa}
En las aplicaciones reales, diseñar la función de recompensa correcta suele ser más difícil que resolver el propio problema de RL:
\begin{itemize}
    \item \textbf{Recompensa dispersa}: en la mayoría de los pasos de tiempo no hay retroalimentación significativa
    \item \textbf{Recompensa mal especificada}: la recompensa diseñada puede producir conductas no previstas (reward hacking)
    \item \textbf{Conflicto entre objetivos múltiples}: es difícil ajustar los pesos entre varios objetivos de recompensa
\end{itemize}
\end{importantbox}

\begin{warningbox}{Reward Hacking}
Cuando la función de recompensa es imperfecta, el agente de RL encuentra formas «legítimas pero no previstas» de maximizar la recompensa. Casos clásicos:
\begin{itemize}
    \item En un juego de carreras, recolectar monedas de recompensa en lugar de completar la pista
    \item Un robot que maximiza la recompensa de velocidad vibrando en lugar de caminar
    \item Un LLM que maximiza la recompensa por longitud repitiendo su salida
\end{itemize}
El reward hacking es uno de los desafíos centrales del problema de alineación (alignment) en RL.
\end{warningbox}

\begin{figure}[H]
\centering
\includegraphics[width=0.9\textwidth]{slides-images/slide-05.jpg}
\caption{El curso ilustra con chatbots, robots que doblan ropa y sistemas de recomendación que, en cuanto una tarea carece de una recompensa confiable o de una métrica verificable objetivamente, el propio planteamiento del problema de RL se vuelve un tema de investigación.}
\end{figure}

\subsubsection{Diseño de los espacios de estados y de acciones}

Elegir una representación de estados y un espacio de acciones adecuados también es una cuestión de diseño importante:
\begin{itemize}
    \item El estado debe contener información suficiente (cumplir la propiedad de Markov), pero sin tener una dimensión demasiado alta
    \item La granularidad del espacio de acciones afecta la eficiencia del aprendizaje
    \item En los LLM, si el espacio de acciones es a nivel de token o a nivel de sentence
\end{itemize}

\subsubsection{Diseño de entornos y puntos de referencia}

\begin{knowledgebox}{Criterios de un buen punto de referencia de RL}
Un punto de referencia de RL útil debe:
\begin{itemize}
    \item Capturar los desafíos centrales de las aplicaciones reales
    \item Permitir iterar con rapidez (con un costo computacional moderado)
    \item Tener métricas de evaluación significativas y confiables
    \item Correlacionarse con el desempeño en el mundo real
\end{itemize}
El desafío actual es que muchos puntos de referencia populares de RL están desconectados de las necesidades de las aplicaciones reales.
\end{knowledgebox}

\subsection{Desafíos en los métodos}

\subsubsection{Eficiencia de muestreo}

Pese a las numerosas mejoras, la eficiencia de muestreo del RL profundo sigue siendo muy inferior a la del aprendizaje humano.

\subsubsection{Capacidad de generalización}

\begin{importantbox}{El problema de la generalización en RL}
La capacidad de generalización de los métodos de RL actuales es limitada:
\begin{itemize}
    \item El desempeño cae abruptamente en entornos nuevos fuera del entorno de entrenamiento
    \item Son excesivamente sensibles a cambios mínimos en la función de recompensa
    \item Les cuesta transferir lo aprendido en una tarea a tareas relacionadas
\end{itemize}
Mejorar la capacidad de generalización del RL es indispensable para desplegarlo en el mundo real.
\end{importantbox}

\subsubsection{Sistemas multiagente y juegos}

El RL multiagente enfrenta desafíos adicionales:
\begin{itemize}
    \item El entorno no es estacionario (las políticas de los demás agentes cambian)
    \item El equilibrio entre coordinación y competencia
    \item El aprendizaje de protocolos de comunicación
\end{itemize}

\begin{figure}[H]
\centering
\includegraphics[width=0.9\textwidth]{slides-images/slide-12.jpg}
\caption{La ponente sitúa la video generation / world model entre las fronteras de los métodos: aportan priors ricos, pero cuando se usan fuera de la distribución de la política se distorsionan por errores físicos y por acciones fuera de distribución.}
\end{figure}

\subsection{Desafíos en el despliegue y la evaluación}

\subsubsection{Seguridad y confiabilidad}

\begin{warningbox}{Riesgos de desplegar sistemas de RL}
Desplegar sistemas de RL en el mundo real implica múltiples riesgos:
\begin{itemize}
    \item Los estados fuera de distribución (out-of-distribution) pueden provocar conductas catastróficas
    \item El proceso de decisión de una política de RL es difícil de explicar
    \item Los ataques adversarios pueden aprovechar las debilidades de la política
    \item Los efectos a largo plazo son difíciles de predecir
\end{itemize}
\end{warningbox}

\subsubsection{Reproducibilidad}

La reproducibilidad de los experimentos de RL profundo ha sido siempre motivo de preocupación:
\begin{itemize}
    \item Sensibilidad a la semilla aleatoria
    \item Alto costo del ajuste de hiperparámetros
    \item Las diferencias entre distintas implementaciones pueden ser muy grandes
\end{itemize}

\subsection{Resumen de la sección}

El RL profundo enfrenta desafíos en todos los niveles, desde la definición del problema hasta los métodos y el despliegue. El diseño de la recompensa, la capacidad de generalización y la seguridad son los problemas abiertos más importantes.
